% Options for packages loaded elsewhere
\PassOptionsToPackage{unicode}{hyperref}
\PassOptionsToPackage{hyphens}{url}
\documentclass[
]{article}
\usepackage{xcolor}
\usepackage{amsmath,amssymb}
\setcounter{secnumdepth}{-\maxdimen} % remove section numbering
\usepackage{iftex}
\ifPDFTeX
  \usepackage[T1]{fontenc}
  \usepackage[utf8]{inputenc}
  \usepackage{textcomp} % provide euro and other symbols
\else % if luatex or xetex
  \usepackage{unicode-math} % this also loads fontspec
  \defaultfontfeatures{Scale=MatchLowercase}
  \defaultfontfeatures[\rmfamily]{Ligatures=TeX,Scale=1}
\fi
\usepackage{lmodern}
\ifPDFTeX\else
  % xetex/luatex font selection
\fi
% Use upquote if available, for straight quotes in verbatim environments
\IfFileExists{upquote.sty}{\usepackage{upquote}}{}
\IfFileExists{microtype.sty}{% use microtype if available
  \usepackage[]{microtype}
  \UseMicrotypeSet[protrusion]{basicmath} % disable protrusion for tt fonts
}{}
\makeatletter
\@ifundefined{KOMAClassName}{% if non-KOMA class
  \IfFileExists{parskip.sty}{%
    \usepackage{parskip}
  }{% else
    \setlength{\parindent}{0pt}
    \setlength{\parskip}{6pt plus 2pt minus 1pt}}
}{% if KOMA class
  \KOMAoptions{parskip=half}}
\makeatother
\usepackage{color}
\usepackage{fancyvrb}
\newcommand{\VerbBar}{|}
\newcommand{\VERB}{\Verb[commandchars=\\\{\}]}
\DefineVerbatimEnvironment{Highlighting}{Verbatim}{commandchars=\\\{\}}
% Add ',fontsize=\small' for more characters per line
\newenvironment{Shaded}{}{}
\newcommand{\AlertTok}[1]{\textcolor[rgb]{1.00,0.00,0.00}{\textbf{#1}}}
\newcommand{\AnnotationTok}[1]{\textcolor[rgb]{0.38,0.63,0.69}{\textbf{\textit{#1}}}}
\newcommand{\AttributeTok}[1]{\textcolor[rgb]{0.49,0.56,0.16}{#1}}
\newcommand{\BaseNTok}[1]{\textcolor[rgb]{0.25,0.63,0.44}{#1}}
\newcommand{\BuiltInTok}[1]{\textcolor[rgb]{0.00,0.50,0.00}{#1}}
\newcommand{\CharTok}[1]{\textcolor[rgb]{0.25,0.44,0.63}{#1}}
\newcommand{\CommentTok}[1]{\textcolor[rgb]{0.38,0.63,0.69}{\textit{#1}}}
\newcommand{\CommentVarTok}[1]{\textcolor[rgb]{0.38,0.63,0.69}{\textbf{\textit{#1}}}}
\newcommand{\ConstantTok}[1]{\textcolor[rgb]{0.53,0.00,0.00}{#1}}
\newcommand{\ControlFlowTok}[1]{\textcolor[rgb]{0.00,0.44,0.13}{\textbf{#1}}}
\newcommand{\DataTypeTok}[1]{\textcolor[rgb]{0.56,0.13,0.00}{#1}}
\newcommand{\DecValTok}[1]{\textcolor[rgb]{0.25,0.63,0.44}{#1}}
\newcommand{\DocumentationTok}[1]{\textcolor[rgb]{0.73,0.13,0.13}{\textit{#1}}}
\newcommand{\ErrorTok}[1]{\textcolor[rgb]{1.00,0.00,0.00}{\textbf{#1}}}
\newcommand{\ExtensionTok}[1]{#1}
\newcommand{\FloatTok}[1]{\textcolor[rgb]{0.25,0.63,0.44}{#1}}
\newcommand{\FunctionTok}[1]{\textcolor[rgb]{0.02,0.16,0.49}{#1}}
\newcommand{\ImportTok}[1]{\textcolor[rgb]{0.00,0.50,0.00}{\textbf{#1}}}
\newcommand{\InformationTok}[1]{\textcolor[rgb]{0.38,0.63,0.69}{\textbf{\textit{#1}}}}
\newcommand{\KeywordTok}[1]{\textcolor[rgb]{0.00,0.44,0.13}{\textbf{#1}}}
\newcommand{\NormalTok}[1]{#1}
\newcommand{\OperatorTok}[1]{\textcolor[rgb]{0.40,0.40,0.40}{#1}}
\newcommand{\OtherTok}[1]{\textcolor[rgb]{0.00,0.44,0.13}{#1}}
\newcommand{\PreprocessorTok}[1]{\textcolor[rgb]{0.74,0.48,0.00}{#1}}
\newcommand{\RegionMarkerTok}[1]{#1}
\newcommand{\SpecialCharTok}[1]{\textcolor[rgb]{0.25,0.44,0.63}{#1}}
\newcommand{\SpecialStringTok}[1]{\textcolor[rgb]{0.73,0.40,0.53}{#1}}
\newcommand{\StringTok}[1]{\textcolor[rgb]{0.25,0.44,0.63}{#1}}
\newcommand{\VariableTok}[1]{\textcolor[rgb]{0.10,0.09,0.49}{#1}}
\newcommand{\VerbatimStringTok}[1]{\textcolor[rgb]{0.25,0.44,0.63}{#1}}
\newcommand{\WarningTok}[1]{\textcolor[rgb]{0.38,0.63,0.69}{\textbf{\textit{#1}}}}
\usepackage{longtable,booktabs,array}
\newcounter{none} % for unnumbered tables
\usepackage{calc} % for calculating minipage widths
% Correct order of tables after \paragraph or \subparagraph
\usepackage{etoolbox}
\makeatletter
\patchcmd\longtable{\par}{\if@noskipsec\mbox{}\fi\par}{}{}
\makeatother
% Allow footnotes in longtable head/foot
\IfFileExists{footnotehyper.sty}{\usepackage{footnotehyper}}{\usepackage{footnote}}
\makesavenoteenv{longtable}
\setlength{\emergencystretch}{3em} % prevent overfull lines
\providecommand{\tightlist}{%
  \setlength{\itemsep}{0pt}\setlength{\parskip}{0pt}}
\usepackage{bookmark}
\IfFileExists{xurl.sty}{\usepackage{xurl}}{} % add URL line breaks if available
\urlstyle{same}
\hypersetup{
  hidelinks,
  pdfcreator={LaTeX via pandoc}}

\author{}
\date{}

\begin{document}

\section{Mesos: A Platform for Fine-Grained Resource Sharing in the Data
Center}\label{mesos-a-platform-for-fine-grained-resource-sharing-in-the-data-center}

Benjamin Hindman, Andy Konwinski, Matei Zaharia, Ali Ghodsi, Anthony D.
Joseph, Randy Katz, Scott Shenker, Ion Stoica University of California,
Berkeley

\section{Abstract}\label{abstract}

We present Mesos, a platform for sharing commodity clusters between
multiple diverse cluster computing frameworks, such as Hadoop and MPI.
Sharing improves cluster utilization and avoids per-framework data
replication. Mesos shares resources in a fine-grained manner, allowing
frameworks to achieve data locality by taking turns reading data stored
on each machine. To support the sophisticated schedulers of today's
frameworks, Mesos introduces a distributed two-level scheduling
mechanism called resource offers. Mesos decides how many resources to
offer each framework, while frameworks decide which resources to accept
and which computations to run on them. Our results show that Mesos can
achieve near-optimal data locality when sharing the cluster among
diverse frameworks, can scale to 50,000 (emulated) nodes, and is
resilient to failures.

\section{1 Introduction}\label{introduction}

Clusters of commodity servers have become a major computing platform,
powering both large Internet services and a growing number of
data-intensive scientific applications. Driven by these applications,
researchers and practitioners have been developing a diverse array of
cluster computing frameworks to simplify programming the cluster.
Prominent examples include MapReduce {[}18{]}, Dryad {[}24{]}, MapReduce
Online {[}17{]} (which supports streaming jobs), Pregel {[}28{]} (a
specialized framework for graph computations), and others {[}27, 19,
30{]}.

It seems clear that new cluster computing frameworks \(^{1}\) will
continue to emerge, and that no framework will be optimal for all
applications. Therefore, organizations will want to run multiple
frameworks in the same cluster, picking the best one for each
application. Multiplexing a cluster between frameworks improves
utilization and allows applications to share access to large datasets
that may be too costly to replicate across clusters.

Two common solutions for sharing a cluster today are either to
statically partition the cluster and run one framework per partition, or
to allocate a set of VMs to each framework. Unfortunately, these
solutions achieve neither high utilization nor efficient data sharing.
The main problem is the mismatch between the allocation granularities of
these solutions and of existing frameworks. Many frameworks, such as
Hadoop and Dryad, employ a fine-grained resource sharing model, where
nodes are subdivided into ``slots'' and jobs are composed of short tasks
that are matched to slots {[}25, 38{]}. The short duration of tasks and
the ability to run multiple tasks per node allow jobs to achieve high
data locality, as each job will quickly get a chance to run on nodes
storing its input data. Short tasks also allow frameworks to achieve
high utilization, as jobs can rapidly scale when new nodes become
available. Unfortunately, because these frameworks are developed
independently, there is no way to perform fine-grained sharing across
frameworks, making it difficult to share clusters and data efficiently
between them.

In this paper, we propose Mesos, a thin resource sharing layer that
enables fine-grained sharing across diverse cluster computing
frameworks, by giving frameworks a common interface for accessing
cluster resources.

The main design question for Mesos is how to build a scalable and
efficient system that supports a wide array of both current and future
frameworks. This is challenging for several reasons. First, each
framework will have different scheduling needs, based on its programming
model, communication pattern, task dependencies, and data placement.
Second, the scheduling system must scale to clusters of tens of
thousands of nodes running hundreds of jobs with millions of tasks.
Finally, because all the applications in the cluster depend on Mesos,
the system must be fault-tolerant and highly available.

One approach would be for Mesos to implement a centralized scheduler
that takes as input framework requirements, resource availability, and
organizational policies, and computes a global schedule for all tasks.
While this

approach can optimize scheduling across frameworks, it faces several
challenges. The first is complexity. The scheduler would need to provide
a sufficiently expressive API to capture all frameworks' requirements,
and to solve an online optimization problem for millions of tasks. Even
if such a scheduler were feasible, this complexity would have a negative
impact on its scalability and resilience. Second, as new frameworks and
new scheduling policies for current frameworks are constantly being
developed {[}37, 38, 40, 26{]}, it is not clear whether we are even at
the point to have a full specification of framework requirements. Third,
many existing frameworks implement their own sophisticated scheduling
{[}25, 38{]}, and moving this functionality to a global scheduler would
require expensive refactoring.

Instead, Mesos takes a different approach: delegating control over
scheduling to the frameworks. This is accomplished through a new
abstraction, called a resource offer, which encapsulates a bundle of
resources that a framework can allocate on a cluster node to run tasks.
Mesos decides how many resources to offer each framework, based on an
organizational policy such as fair sharing, while frameworks decide
which resources to accept and which tasks to run on them. While this
decentralized scheduling model may not always lead to globally optimal
scheduling, we have found that it performs surprisingly well in
practice, allowing frameworks to meet goals such as data locality nearly
perfectly. In addition, resource offers are simple and efficient to
implement, allowing Mesos to be highly scalable and robust to failures.

Mesos also provides other benefits to practitioners. First, even
organizations that only use one framework can use Mesos to run multiple
instances of that framework in the same cluster, or multiple versions of
the framework. Our contacts at Yahoo! and Facebook indicate that this
would be a compelling way to isolate production and experimental Hadoop
workloads and to roll out new versions of Hadoop {[}11, 10{]}. Second,
Mesos makes it easier to develop and immediately experiment with new
frameworks. The ability to share resources across multiple frameworks
frees the developers to build and run specialized frameworks targeted at
particular problem domains rather than one-size-fits-all abstractions.
Frameworks can therefore evolve faster and provide better support for
each problem domain.

We have implemented Mesos in 10,000 lines of C++. The system scales to
50,000 (emulated) nodes and uses ZooKeeper {[}4{]} for fault tolerance.
To evaluate Mesos, we have ported three cluster computing systems to run
over it: Hadoop, MPI, and the Torque batch scheduler. To validate our
hypothesis that specialized frameworks provide value over general ones,
we have also built a new framework on top of Mesos called Spark,
optimized for iterative jobs where a dataset is reused in many parallel
operations, and shown that Spark can outperform Hadoop by 10x in
iterative machine learning workloads.

\textit{[Image omitted: images/7eaaa664d5ce4baa22872dbe26492a6c1f45d3dd276abf593f26a58438d86e0e.jpg]}

line

{\def\LTcaptype{none} % do not increment counter
\begin{longtable}[]{@{}lll@{}}
\toprule\noalign{}
Duration (s) & MapReduce Jobs & Map \& Reduce Tasks \\
\midrule\noalign{}
\endhead
\bottomrule\noalign{}
\endlastfoot
1 & 0.0 & 0.0 \\
10 & 0.1 & 0.3 \\
100 & 0.5 & 0.7 \\
1000 & 0.9 & 0.95 \\
10000 & 0.95 & 0.98 \\
100000 & 0.98 & 0.99 \\
\end{longtable}
}

Figure 1: CDF of job and task durations in Facebook's Hadoop data
warehouse (data from {[}38{]}).

This paper is organized as follows. Section 2 details the data center
environment that Mesos is designed for. Section 3 presents the
architecture of Mesos. Section 4 analyzes our distributed scheduling
model (resource offers) and characterizes the environments that it works
well in. We present our implementation of Mesos in Section 5 and
evaluate it in Section 6. We survey related work in Section 7. Finally,
we conclude in Section 8.

\section{2 Target Environment}\label{target-environment}

As an example of a workload we aim to support, consider the Hadoop data
warehouse at Facebook {[}5{]}. Facebook loads logs from its web services
into a 2000-node Hadoop cluster, where they are used for applications
such as business intelligence, spam detection, and ad optimization. In
addition to ``production'' jobs that run periodically, the cluster is
used for many experimental jobs, ranging from multi-hour machine
learning computations to 1-2 minute ad-hoc queries submitted
interactively through an SQL interface called Hive {[}3{]}. Most jobs
are short (the median job being 84s long), and the jobs are composed of
fine-grained map and reduce tasks (the median task being 23s), as shown
in Figure 1.

To meet the performance requirements of these jobs, Facebook uses a fair
scheduler for Hadoop that takes advantage of the fine-grained nature of
the workload to allocate resources at the level of tasks and to optimize
data locality {[}38{]}. Unfortunately, this means that the cluster can
only run Hadoop jobs. If a user wishes to write an ad targeting
algorithm in MPI instead of MapReduce, perhaps because MPI is more
efficient for this job's communication pattern, then the user must set
up a separate MPI cluster and import terabytes of data into it. This
problem is not hypothetical; our contacts at Yahoo! and Facebook report
that users want to run MPI and MapReduce Online (a streaming MapReduce)
{[}11, 10{]}. Mesos aims to provide fine-grained sharing between
multiple cluster computing frameworks to enable these usage scenarios.

\textit{[Image omitted: images/92d35367fff062f628f54e7257600cc38972e74b18aa9fd1aaf9604a8989ef82.jpg]}

flowchart

\begin{Shaded}
\begin{Highlighting}[]
\NormalTok{graph TD}
\NormalTok{    A["Hadoop scheduler"] {-}{-}\textgreater{} C["Mesos master"]}
\NormalTok{    B["MPI scheduler"] {-}{-}\textgreater{} C}
\NormalTok{    C {-}{-}\textgreater{} D["Mesos slave"]}
\NormalTok{    C {-}{-}\textgreater{} E["Mesos slave"]}
\NormalTok{    C {-}{-}\textgreater{} F["Mesos slave"]}
\NormalTok{    D {-}{-}\textgreater{} G["Hadoop executor task"]}
\NormalTok{    E {-}{-}\textgreater{} H["MPI executor task"]}
\NormalTok{    F {-}{-}\textgreater{} I["Hadoop executor task"]}
\NormalTok{    F {-}{-}\textgreater{} J["MPI executor task"]}
\NormalTok{    C {-}.{-}\textgreater{} K["ZooKeeper quorum"]}
\NormalTok{    K {-}{-}\textgreater{} L["Standby master"]}
\NormalTok{    K {-}{-}\textgreater{} M["Standby master"]}
\end{Highlighting}
\end{Shaded}

Figure 2: Mesos architecture diagram, showing two running frameworks
(Hadoop and MPI).

\section{3 Architecture}\label{architecture}

We begin our description of Mesos by discussing our design philosophy.
We then describe the components of Mesos, our resource allocation
mechanisms, and how Mesos achieves isolation, scalability, and fault
tolerance.

\section{3.1 Design Philosophy}\label{design-philosophy}

Mesos aims to provide a scalable and resilient core for enabling various
frameworks to efficiently share clusters. Because cluster frameworks are
both highly diverse and rapidly evolving, our overriding design
philosophy has been to define a minimal interface that enables efficient
resource sharing across frameworks, and otherwise push control of task
scheduling and execution to the frameworks. Pushing control to the
frameworks has two benefits. First, it allows frameworks to implement
diverse approaches to various problems in the cluster (e.g., achieving
data locality, dealing with faults), and to evolve these solutions
independently. Second, it keeps Mesos simple and minimizes the rate of
change required of the system, which makes it easier to keep Mesos
scalable and robust.

Although Mesos provides a low-level interface, we expect higher-level
libraries implementing common functionality (such as fault tolerance) to
be built on top of it. These libraries would be analogous to library
OSes in the exokernel {[}20{]}. Putting this functionality in libraries
rather than in Mesos allows Mesos to remain small and flexible, and lets
the libraries evolve independently.

\section{3.2 Overview}\label{overview}

Figure 2 shows the main components of Mesos. Mesos consists of a master
process that manages slave daemons running on each cluster node, and
frameworks that run tasks on these slaves.

The master implements fine-grained sharing across frameworks using
resource offers. Each resource offer is a list of free resources on
multiple slaves. The master decides how many resources to offer to each
framework according to an organizational policy, such as fair sharing

\textit{[Image omitted: images/1905bdee4853014917d935d612a86876ba3ac274011d3c7dfb05a592c7694fc4.jpg]}

flowchart

\begin{Shaded}
\begin{Highlighting}[]
\NormalTok{graph TD}
\NormalTok{    A["Framework 1\textbackslash{}nJob 1\textbackslash{}nJob 2\textbackslash{}nFW Scheduler"] {-}{-}\textgreater{}|2| B["Allocation module"]}
\NormalTok{    C["Framework 2\textbackslash{}nJob 1\textbackslash{}nJob 2\textbackslash{}nFW Scheduler"] {-}{-}\textgreater{}|3| B}
\NormalTok{    B {-}{-}\textgreater{}|4| D["Slave 1\textbackslash{}nExecutor\textbackslash{}nTask\textbackslash{}nTask"]}
\NormalTok{    B {-}{-}\textgreater{}|1| D}
\NormalTok{    B {-}{-}\textgreater{}|4| D}
\NormalTok{    B {-}{-}\textgreater{}|4| E["Mesos master"]}
\NormalTok{    E {-}{-}\textgreater{}|4| D}
\NormalTok{    style A fill:\#f9f,stroke:\#333}
\NormalTok{    style C fill:\#f9f,stroke:\#333}
\NormalTok{    style B fill:\#ccf,stroke:\#333}
\NormalTok{    style D fill:\#cfc,stroke:\#333}
\NormalTok{    style E fill:\#fcc,stroke:\#333}
\end{Highlighting}
\end{Shaded}

Figure 3: Resource offer example.

or priority. To support a diverse set of inter-framework allocation
policies, Mesos lets organizations define their own policies via a
pluggable allocation module.

Each framework running on Mesos consists of two components: a scheduler
that registers with the master to be offered resources, and an executor
process that is launched on slave nodes to run the framework's tasks.
While the master determines how many resources to offer to each
framework, the frameworks' schedulers select which of the offered
resources to use. When a framework accepts offered resources, it passes
Mesos a description of the tasks it wants to launch on them.

Figure 3 shows an example of how a framework gets scheduled to run
tasks. In step (1), slave 1 reports to the master that it has 4 CPUs and
4 GB of memory free. The master then invokes the allocation module,
which tells it that framework 1 should be offered all available
resources. In step (2), the master sends a resource offer describing
these resources to framework 1. In step (3), the framework's scheduler
replies to the master with information about two tasks to run on the
slave, using \(\langle2\) CPUs, 1 GB RAM \(\rangle\) for the first task,
and \(\langle1\) CPUs, 2 GB RAM \(\rangle\) for the second task.
Finally, in step (4), the master sends the tasks to the slave, which
allocates appropriate resources to the framework's executor, which in
turn launches the two tasks (depicted with dotted borders). Because 1
CPU and 1 GB of RAM are still free, the allocation module may now offer
them to framework 2. In addition, this resource offer process repeats
when tasks finish and new resources become free.

To maintain a thin interface and enable frameworks to evolve
independently, Mesos does not require frameworks to specify their
resource requirements or constraints. Instead, Mesos gives frameworks
the ability to reject offers. A framework can reject resources that do
not satisfy its constraints in order to wait for ones that do. Thus, the
rejection mechanism enables frameworks to support arbitrarily complex
resource constraints while keeping Mesos simple and scalable.

One potential challenge with solely using the reject-

tion mechanism to satisfy all framework constraints is efficiency: a
framework may have to wait a long time before it receives an offer
satisfying its constraints, and Mesos may have to send an offer to many
frameworks before one of them accepts it. To avoid this, Mesos also
allows frameworks to set filters, which are Boolean predicates
specifying that a framework will always reject certain resources. For
example, a framework might specify a whitelist of nodes it can run on.

There are two points worth noting. First, filters represent just a
performance optimization for the resource offer model, as the frameworks
still have the ultimate control to reject any resources that they cannot
express filters for and to choose which tasks to run on each node.
Second, as we will show in this paper, when the workload consists of
fine-grained tasks (e.g., in MapReduce and Dryad workloads), the
resource offer model performs surprisingly well even in the absence of
filters. In particular, we have found that a simple policy called delay
scheduling {[}38{]}, in which frameworks wait for a limited time to
acquire nodes storing their data, yields nearly optimal data locality
with a wait time of 1-5s.

In the rest of this section, we describe how Mesos performs two key
functions: resource allocation (§3.3) and resource isolation (§3.4). We
then describe filters and several other mechanisms that make resource
offers scalable and robust (§3.5). Finally, we discuss fault tolerance
in Mesos (§3.6) and summarize the Mesos API (§3.7).

\section{3.3 Resource Allocation}\label{resource-allocation}

Mesos delegates allocation decisions to a pluggable allocation module,
so that organizations can tailor allocation to their needs. So far, we
have implemented two allocation modules: one that performs fair sharing
based on a generalization of max-min fairness for multiple resources
{[}21{]} and one that implements strict priorities. Similar policies are
used in Hadoop and Dryad {[}25, 38{]}.

In normal operation, Mesos takes advantage of the fact that most tasks
are short, and only reallocates resources when tasks finish. This
usually happens frequently enough so that new frameworks acquire their
share quickly. For example, if a framework's share is 10\% of the
cluster, it needs to wait approximately 10\% of the mean task length to
receive its share. However, if a cluster becomes filled by long tasks,
e.g., due to a buggy job or a greedy framework, the allocation module
can also revoke (kill) tasks. Before killing a task, Mesos gives its
framework a grace period to clean it up.

We leave it up to the allocation module to select the policy for
revoking tasks, but describe two related mechanisms here. First, while
killing a task has a low impact on many frameworks (e.g., MapReduce), it
is harmful for frameworks with interdependent tasks (e.g., MPI). We
allow these frameworks to avoid being killed by letting allocation
modules expose a guaranteed allocation to each framework---a quantity of
resources that the framework may hold without losing tasks. Frameworks
read their guaranteed allocations through an API call. Allocation
modules are responsible for ensuring that the guaranteed allocations
they provide can all be met concurrently. For now, we have kept the
semantics of guaranteed allocations simple: if a framework is below its
guaranteed allocation, none of its tasks should be killed, and if it is
above, any of its tasks may be killed.

Second, to decide when to trigger revocation, Mesos must know which of
the connected frameworks would use more resources if they were offered
them. Frameworks indicate their interest in offers through an API call.

\section{3.4 Isolation}\label{isolation}

Mesos provides performance isolation between framework executors running
on the same slave by leveraging existing OS isolation mechanisms. Since
these mechanisms are platform-dependent, we support multiple isolation
mechanisms through pluggable isolation modules.

We currently isolate resources using OS container technologies,
specifically Linux Containers {[}9{]} and Solaris Projects {[}13{]}.
These technologies can limit the CPU, memory, network bandwidth, and (in
new Linux kernels) I/O usage of a process tree. These isolation
technologies are not perfect, but using containers is already an
advantage over frameworks like Hadoop, where tasks from different jobs
simply run in separate processes.

\section{3.5 Making Resource Offers Scalable and
Robust}\label{making-resource-offers-scalable-and-robust}

Because task scheduling in Mesos is a distributed process, it needs to
be efficient and robust to failures. Mesos includes three mechanisms to
help with this goal.

First, because some frameworks will always reject certain resources,
Mesos lets them short-circuit the rejection process and avoid
communication by providing filters to the master. We currently support
two types of filters: ``only offer nodes from list L'' and ``only offer
nodes with at least R resources free''. However, other types of
predicates could also be supported. Note that unlike generic constraint
languages, filters are Boolean predicates that specify whether a
framework will reject one bundle of resources on one node, so they can
be evaluated quickly on the master. Any resource that does not pass a
framework's filter is treated exactly like a rejected resource.

Second, because a framework may take time to respond to an offer, Mesos
counts resources offered to a framework towards its allocation of the
cluster. This is a strong incentive for frameworks to respond to offers
quickly and to filter resources that they cannot use.

Third, if a framework has not responded to an offer for a sufficiently
long time, Mesos rescinds the offer and re-offers the resources to other
frameworks.

\section{Scheduler Callbacks}\label{scheduler-callbacks}

resourceOffer(offerId, offers) offerRescinded(offerId)
statusUpdate(taskId, status) slaveLost(slaveld)

\section{Scheduler Actions}\label{scheduler-actions}

replyToOffer(offerId, tasks) setNeedsOffers(bool) setFilters(filters)
getGuaranteedShare() killTask(taskId)

\section{Executor Callbacks}\label{executor-callbacks}

launchTask(taskDescriptor) killTask(taskId)

\section{Executor Actions}\label{executor-actions}

sendStatus(taskId, status)

Table 1: Mesos API functions for schedulers and executors.

\section{3.6 Fault Tolerance}\label{fault-tolerance}

Since all the frameworks depend on the Mesos master, it is critical to
make the master fault-tolerant. To achieve this, we have designed the
master to be soft state, so that a new master can completely reconstruct
its internal state from information held by the slaves and the framework
schedulers. In particular, the master's only state is the list of active
slaves, active frameworks, and running tasks. This information is
sufficient to compute how many resources each framework is using and run
the allocation policy. We run multiple masters in a hot-standby
configuration using ZooKeeper {[}4{]} for leader election. When the
active master fails, the slaves and schedulers connect to the next
elected master and repopulate its state.

Aside from handling master failures, Mesos reports node failures and
executor crashes to frameworks' schedulers. Frameworks can then react to
these failures using the policies of their choice.

Finally, to deal with scheduler failures, Mesos allows a framework to
register multiple schedulers such that when one fails, another one is
notified by the Mesos master to take over. Frameworks must use their own
mechanisms to share state between their schedulers.

\section{3.7 API Summary}\label{api-summary}

Table 1 summarizes the Mesos API. The ``callback'' columns list
functions that frameworks must implement, while ``actions'' are
operations that they can invoke.

\section{4 Mesos Behavior}\label{mesos-behavior}

In this section, we study Mesos's behavior for different workloads. Our
goal is not to develop an exact model of the system, but to provide a
coarse understanding of its behavior, in order to characterize the
environments that Mesos's distributed scheduling model works well in.

In short, we find that Mesos performs very well when frameworks can
scale up and down elastically, tasks durations are homogeneous, and
frameworks prefer all nodes equally ( \(\S4.2\) ). When different
frameworks prefer different nodes, we show that Mesos can emulate a
centralized scheduler that performs fair sharing across frameworks (
\(\S4.3\) ). In addition, we show that Mesos can handle heterogeneous
task durations without impacting the performance of frameworks with
short tasks (§4.4). We also discuss how frameworks are incentivized to
improve their performance under Mesos, and argue that these incentives
also improve overall cluster utilization (§4.5). We conclude this
section with some limitations of Mesos's distributed scheduling model
(§4.6).

\section{4.1 Definitions, Metrics and
Assumptions}\label{definitions-metrics-and-assumptions}

In our discussion, we consider three metrics:

\begin{itemize}
\tightlist
\item
  Framework ramp-up time: time it takes a new framework to achieve its
  allocation (e.g., fair share);\\
\item
  Job completion time: time it takes a job to complete, assuming one job
  per framework;\\
\item
  System utilization: total cluster utilization.
\end{itemize}

We characterize workloads along two dimensions: elasticity and task
duration distribution. An elastic framework, such as Hadoop and Dryad,
can scale its resources up and down, i.e., it can start using nodes as
soon as it acquires them and release them as soon its task finish. In
contrast, a rigid framework, such as MPI, can start running its jobs
only after it has acquired a fixed quantity of resources, and cannot
scale up dynamically to take advantage of new resources or scale down
without a large impact on performance. For task durations, we consider
both homogeneous and heterogeneous distributions.

We also differentiate between two types of resources: mandatory and
preferred. A resource is mandatory if a framework must acquire it in
order to run. For example, a graphical processing unit (GPU) is
mandatory if a framework cannot run without access to GPU. In contrast,
a resource is preferred if a framework performs ``better'' using it, but
can also run using another equivalent resource. For example, a framework
may prefer running on a node that locally stores its data, but may also
be able to read the data remotely if it must.

We assume the amount of mandatory resources requested by a framework
never exceeds its guaranteed share. This ensures that frameworks will
not deadlock waiting for the mandatory resources to become free.
\(^{2}\) For simplicity, we also assume that all tasks have the same
resource demands and run on identical slices of machines called slots,
and that each framework runs a single job.

\section{4.2 Homogeneous Tasks}\label{homogeneous-tasks}

We consider a cluster with n slots and a framework, f, that is entitled
to k slots. For the purpose of this analysis, we consider two
distributions of the task durations: constant (i.e., all tasks have the
same length) and exponential. Let the mean task duration be T, and
assume that framework f runs a job which requires \(\beta kT\) total
com-

Elastic Framework

Rigid Framework

Constant dist.

Exponential dist.

Constant dist.

Exponential dist.

Ramp-up time

T

T ln k

T

T ln k

Completion time

(1/2 + β)T

(1 + β)T

(1 + β)T

(ln k + β)T

Utilization

1

1

β/(1/2 + β)

β/(ln k - 1 + β)

Table 2: Ramp-up time, job completion time and utilization for both
elastic and rigid frameworks, and for both constant and exponential task
duration distributions. The framework starts with no slots. k is the
number of slots the framework is entitled under the scheduling policy,
and \(\beta T\) represents the time it takes a job to complete assuming
the framework gets all k slots at once.

putation time. That is, when the framework has k slots, it takes its job
\(\beta T\) time to finish.

Table 2 summarizes the job completion times and system utilization for
the two types of frameworks and the two types of task length
distributions. As expected, elastic frameworks with constant task
durations perform the best, while rigid frameworks with exponential task
duration perform the worst. Due to lack of space, we present only the
results here and include derivations in {[}23{]}.

Framework ramp-up time: If task durations are constant, it will take
framework f at most T time to acquire k slots. This is simply because
during a T interval, every slot will become available, which will enable
Mesos to offer the framework all k of its preferred slots. If the
duration distribution is exponential, the expected ramp-up time can be
as high as \(T \ln k\) {[}23{]}.

Job completion time: The expected completion time \(^{3}\) of an elastic
job is at most \((1 + \beta)T\) , which is within T (i.e., the mean task
duration) of the completion time of the job when it gets all its slots
instantaneously. Rigid jobs achieve similar completion times for
constant task durations, but exhibit much higher completion times for
exponential job durations, i.e., \((\ln k + \beta)T\) . This is simply
because it takes a framework \(T \ln k\) time on average to acquire all
its slots and be able to start its job.

System utilization: Elastic jobs fully utilize their allocated slots,
because they can use every slot as soon as they get it. As a result,
assuming infinite demand, a system running only elastic jobs is fully
utilized. Rigid frameworks achieve slightly worse utilizations, as their
jobs cannot start before they get their full allocations, and thus they
waste the resources held while ramping up.

\section{4.3 Placement Preferences}\label{placement-preferences}

So far, we have assumed that frameworks have no slot preferences. In
practice, different frameworks prefer different nodes and their
preferences may change over time. In this section, we consider the case
where frameworks have different preferred slots.

The natural question is how well Mesos will work compared to a central
scheduler that has full information about framework preferences. We
consider two cases: (a) there exists a system configuration in which
each framework gets all its preferred slots and achieves its full
allocation, and (b) there is no such configuration, i.e., the demand for
some preferred slots exceeds the supply.

In the first case, it is easy to see that, irrespective of the initial
configuration, the system will converge to the state where each
framework allocates its preferred slots after at most one T interval.
This is simple because during a T interval all slots become available,
and as a result each framework will be offered its preferred slots.

In the second case, there is no configuration in which all frameworks
can satisfy their preferences. The key question in this case is how
should one allocate the preferred slots across the frameworks demanding
them. In particular, assume there are p slots preferred by m frameworks,
where framework i requests \(r_{i}\) such slots, and
\(\sum_{i=1}^{m} r_{i} > x\) . While many allocation policies are
possible, here we consider a weighted fair allocation policy where the
weight associated with framework i is its intended total allocation,
\(s_{i}\) . In other words, assuming that each framework has enough
demand, we aim to allocate \(p \cdot s_{i} / (\sum_{i=1}^{m} s_{i})\)
preferred slots to framework i.

The challenge in Mesos is that the scheduler does not know the
preferences of each framework. Fortunately, it turns out that there is
an easy way to achieve the weighted allocation of the preferred slots
described above: simply perform lottery scheduling {[}36{]}, offering
slots to frameworks with probabilities proportional to their intended
allocations. In particular, when a slot becomes available, Mesos can
offer that slot to framework i with probability
\(s_{i}/(\sum_{i=1}^{n} s_{i})\) , where n is the total number of
frameworks in the system. Furthermore, because each framework i receives
on average \(s_{i}\) slots every T time units, the results for ramp-up
times and completion times in Section 4.2 still hold.

\section{4.4 Heterogeneous Tasks}\label{heterogeneous-tasks}

So far we have assumed that frameworks have homogeneous task duration
distributions, i.e., that all frameworks have the same task duration
distribution. In this section, we discuss frameworks with heterogeneous
task duration distributions. In particular, we consider a workload where
tasks that are either short and long, where the mean duration of the
long tasks is significantly longer than the mean of the short tasks.
Such heterogeneous

workloads can hurt frameworks with short tasks. In the worst case, all
nodes required by a short job might be filled with long tasks, so the
job may need to wait a long time (relative to its execution time) to
acquire resources.

We note first that random task assignment can work well if the fraction
\(\phi\) of long tasks is not very close to 1 and if each node supports
multiple slots. For example, in a cluster with \(S\) slots per node, the
probability that a node is filled with long tasks will be \(\phi^S\) .
When \(S\) is large (e.g., in the case of multicore machines), this
probability is small even with \(\phi > 0.5\) . If \(S = 8\) and
\(\phi = 0.5\) , for example, the probability that a node is filled with
long tasks is \(0.4\%\) . Thus, a framework with short tasks can still
acquire many preferred slots in a short period of time. In addition, the
more slots a framework is able to use, the likelier it is that at least
\(k\) of them are running short tasks.

To further alleviate the impact of long tasks, Mesos can be extended
slightly to allow allocation policies to reserve some resources on each
node for short tasks. In particular, we can associate a maximum task
duration with some of the resources on each node, after which tasks
running on those resources are killed. These time limits can be exposed
to the frameworks in resource offers, allowing them to choose whether to
use these resources. This scheme is similar to the common policy of
having a separate queue for short jobs in HPC clusters.

\section{4.5 Framework Incentives}\label{framework-incentives}

Mesos implements a decentralized scheduling model, where each framework
decides which offers to accept. As with any decentralized system, it is
important to understand the incentives of entities in the system. In
this section, we discuss the incentives of frameworks (and their users)
to improve the response times of their jobs.

Short tasks: A framework is incentivized to use short tasks for two
reasons. First, it will be able to allocate any resources reserved for
short slots. Second, using small tasks minimizes the wasted work if the
framework loses a task, either due to revocation or simply due to
failures.

Scale elastically: The ability of a framework to use resources as soon
as it acquires them---instead of waiting to reach a given minimum
allocation---would allow the framework to start (and complete) its jobs
earlier. In addition, the ability to scale up and down allows a
framework to grab unused resources opportunistically, as it can later
release them with little negative impact.

Do not accept unknown resources: Frameworks are incentivized not to
accept resources that they cannot use because most allocation policies
will count all the resources that a framework owns when making offers.

We note that these incentives align well with our goal of improving
utilization. If frameworks use short tasks, Mesos can reallocate
resources quickly between them, reducing latency for new jobs and wasted
work for revocation. If frameworks are elastic, they will
opportunistically utilize all the resources they can obtain. Finally, if
frameworks do not accept resources that they do not understand, they
will leave them for frameworks that do.

We also note that these properties are met by many current cluster
computing frameworks, such as MapReduce and Dryad, simply because using
short independent tasks simplifies load balancing and fault recovery.

\section{4.6 Limitations of Distributed
Scheduling}\label{limitations-of-distributed-scheduling}

Although we have shown that distributed scheduling works well in a range
of workloads relevant to current cluster environments, like any
decentralized approach, it can perform worse than a centralized
scheduler. We have identified three limitations of the distributed
model:

Fragmentation: When tasks have heterogeneous resource demands, a
distributed collection of frameworks may not be able to optimize bin
packing as well as a centralized scheduler. However, note that the
wasted space due to suboptimal bin packing is bounded by the ratio
between the largest task size and the node size. Therefore, clusters
running ``larger'' nodes (e.g., multicore nodes) and ``smaller'' tasks
within those nodes will achieve high utilization even with distributed
scheduling.

There is another possible bad outcome if allocation modules reallocate
resources in a naïve manner: when a cluster is filled by tasks with
small resource requirements, a framework f with large resource
requirements may starve, because whenever a small task finishes, f
cannot accept the resources freed by it, but other frameworks can. To
accommodate frameworks with large per-task resource requirements,
allocation modules can support a minimum offer size on each slave, and
abstain from offering resources on the slave until this amount is free.

Interdependent framework constraints: It is possible to construct
scenarios where, because of esoteric interdependencies between
frameworks (e.g., certain tasks from two frameworks cannot be
colocated), only a single global allocation of the cluster performs
well. We argue such scenarios are rare in practice. In the model
discussed in this section, where frameworks only have preferences over
which nodes they use, we showed that allocations approximate those of
optimal schedulers.

Framework complexity: Using resource offers may make framework
scheduling more complex. We argue, however, that this difficulty is not
onerous. First, whether using Mesos or a centralized scheduler,
frameworks need to know their preferences; in a centralized scheduler,
the framework needs to express them to the scheduler, whereas in Mesos,
it must use them to decide which offers to accept. Second, many
scheduling policies for existing frameworks are online algorithms,
because frame-

works cannot predict task times and must be able to handle failures and
stragglers {[}18, 40, 38{]}. These policies are easy to implement over
resource offers.

\section{5 Implementation}\label{implementation}

We have implemented Mesos in about 10,000 lines of C++. The system runs
on Linux, Solaris and OS X, and supports frameworks written in C++,
Java, and Python.

To reduce the complexity of our implementation, we use a C++ library
called libprocess {[}7{]} that provides an actor-based programming model
using efficient asynchronous I/O mechanisms (epoll, kqueue, etc). We
also use ZooKeeper {[}4{]} to perform leader election.

Mesos can use Linux containers {[}9{]} or Solaris projects {[}13{]} to
isolate tasks. We currently isolate CPU cores and memory. We plan to
leverage recently added support for network and I/O isolation in Linux
{[}8{]} in the future.

We have implemented four frameworks on top of Mesos. First, we have
ported three existing cluster computing systems: Hadoop {[}2{]}, the
Torque resource scheduler {[}33{]}, and the MPICH2 implementation of MPI
{[}16{]}. None of these ports required changing these frameworks' APIs,
so all of them can run unmodified user programs. In addition, we built a
specialized framework for iterative jobs called Spark, which we discuss
in Section 5.3.

\section{5.1 Hadoop Port}\label{hadoop-port}

Porting Hadoop to run on Mesos required relatively few modifications,
because Hadoop's fine-grained map and reduce tasks map cleanly to Mesos
tasks. In addition, the Hadoop master, known as the JobTracker, and
Hadoop slaves, known as TaskTrackers, fit naturally into the Mesos model
as a framework scheduler and executor.

To add support for running Hadoop on Mesos, we took advantage of the
fact that Hadoop already has a pluggable API for writing job schedulers.
We wrote a Hadoop scheduler that connects to Mesos, launches
TaskTrackers as its executors, and maps each Hadoop task to a Mesos
task. When there are unlaunched tasks in Hadoop, our scheduler first
starts Mesos tasks on the nodes of the cluster that it wants to use, and
then sends the Hadoop tasks to them using Hadoop's existing internal
interfaces. When tasks finish, our executor notifies Mesos by listening
for task finish events using an API in the TaskTracker.

We used delay scheduling {[}38{]} to achieve data locality by waiting
for slots on the nodes that contain task input data. In addition, our
approach allowed us to reuse Hadoop's existing logic for re-scheduling
of failed tasks and for speculative execution (straggler mitigation).

We also needed to change how map output data is served to reduce tasks.
Hadoop normally writes map output files to the local filesystem, then
serves these to reduce tasks using an HTTP server included in the
TaskTracker. However, the TaskTracker within Mesos runs as an executor,
which may be terminated if it is not running tasks. This would make map
output files unavailable to reduce tasks. We solved this problem by
providing a shared file server on each node in the cluster to serve
local files. Such a service is useful beyond Hadoop, to other frameworks
that write data locally on each node.

In total, our Hadoop port is 1500 lines of code.

\section{5.2 Torque and MPI Ports}\label{torque-and-mpi-ports}

We have ported the Torque cluster resource manager to run as a framework
on Mesos. The framework consists of a Mesos scheduler and executor,
written in 360 lines of Python code, that launch and manage different
components of Torque. In addition, we modified 3 lines of Torque source
code to allow it to elastically scale up and down on Mesos depending on
the jobs in its queue.

After registering with the Mesos master, the framework scheduler
configures and launches a Torque server and then periodically monitors
the server's job queue. While the queue is empty, the scheduler releases
all tasks (down to an optional minimum, which we set to 0) and refuses
all resource offers it receives from Mesos. Once a job gets added to
Torque's queue (using the standard qsub command), the scheduler begins
accepting new resource offers. As long as there are jobs in Torque's
queue, the scheduler accepts offers as necessary to satisfy the
constraints of as many jobs in the queue as possible. On each node where
offers are accepted, Mesos launches our executor, which in turn starts a
Torque backend daemon and registers it with the Torque server. When
enough Torque backend daemons have registered, the torque server will
launch the next job in its queue.

Because jobs that run on Torque (e.g.~MPI) may not be fault tolerant,
Torque avoids having its tasks revoked by not accepting resources beyond
its guaranteed allocation.

In addition to the Torque framework, we also created a Mesos MPI
``wrapper'' framework, written in 200 lines of Python code, for running
MPI jobs directly on Mesos.

\section{5.3 Spark Framework}\label{spark-framework}

Mesos enables the creation of specialized frameworks optimized for
workloads for which more general execution layers may not be optimal. To
test the hypothesis that simple specialized frameworks provide value, we
identified one class of jobs that were found to perform poorly on Hadoop
by machine learning researchers at our lab: iterative jobs, where a
dataset is reused across a number of iterations. We built a specialized
framework called Spark {[}39{]} optimized for these workloads.

One example of an iterative algorithm used in machine learning is
logistic regression {[}22{]}. This algorithm seeks to find a line that
separates two sets of labeled data points. The algorithm starts with a
random line w. Then, on each iteration, it computes the gradient of an
objective

\textit{[Image omitted: images/c3313ead8321b507d91d34f5b32750b319236320d33e39c4c35b7a88e368c3b5.jpg]}

flowchart

\begin{Shaded}
\begin{Highlighting}[]
\NormalTok{graph TD}
\NormalTok{    subgraph Input Layer}
\NormalTok{        A["Input Layer"] {-}{-}\textgreater{}|x1| B["Weight"]}
\NormalTok{        A {-}{-}\textgreater{}|x1| C["Activation Function"]}
\NormalTok{        B {-}{-}\textgreater{} D["Weight"]}
\NormalTok{        C {-}{-}\textgreater{} D}
\NormalTok{        D {-}{-}\textgreater{} E["f(x,w)"]}
\NormalTok{    end}
\NormalTok{    subgraph Output Layer}
\NormalTok{        F["Input Layer"] {-}{-}\textgreater{}|x1| G["Weight"]}
\NormalTok{        F {-}{-}\textgreater{}|x1| H["Activation Function"]}
\NormalTok{        G {-}{-}\textgreater{} I["Weight"]}
\NormalTok{        H {-}{-}\textgreater{} I}
\NormalTok{        I {-}{-}\textgreater{} J["f(x,w)"]}
\NormalTok{    end}
\NormalTok{    style Input Layer fill:\#f9f,stroke:\#333}
\NormalTok{    style Output Layer fill:\#bbf,stroke:\#333}
\end{Highlighting}
\end{Shaded}

\begin{enumerate}
\def\labelenumi{\alph{enumi})}
\tightlist
\item
  Dryad
\end{enumerate}

\textit{[Image omitted: images/6dd9bf4e01d56a29fdf19caf446ec13bf5ad8af8e4af8448ceed6fb39feffc9b.jpg]}

flowchart

\begin{Shaded}
\begin{Highlighting}[]
\NormalTok{graph TD}
\NormalTok{    A["f(x,w)"] {-}{-}\textgreater{} B(( ))}
\NormalTok{    B {-}{-}\textgreater{} C["w"]}
\NormalTok{    B {-}{-}\textgreater{} D["x"]}
\NormalTok{    C {-}{-}\textgreater{} E[" "]}
\NormalTok{    D {-}{-}\textgreater{} F[" "]}
\NormalTok{    E {-}{-}\textgreater{} G[" "]}
\NormalTok{    F {-}{-}\textgreater{} H[" "]}
\NormalTok{    G {-}{-}\textgreater{} I[" "]}
\NormalTok{    H {-}{-}\textgreater{} I}
\NormalTok{    I {-}{-}\textgreater{} J[" "]}
\NormalTok{    style A fill:\#f9f,stroke:\#333}
\NormalTok{    style J fill:\#ccf,stroke:\#333}
\end{Highlighting}
\end{Shaded}

\begin{enumerate}
\def\labelenumi{\alph{enumi})}
\setcounter{enumi}{1}
\tightlist
\item
  Spark\\
  Figure 4: Data flow of a logistic regression job in Dryad vs.~Spark.
  Solid lines show data flow within the framework. Dashed lines show
  reads from a distributed file system. Spark reuses in-memory data
  across iterations to improve efficiency.
\end{enumerate}

function that measures how well the line separates the points, and
shifts w along this gradient. This gradient computation amounts to
evaluating a function \(f(x,w)\) over each data point x and summing the
results. An implementation of logistic regression in Hadoop must run
each iteration as a separate MapReduce job, because each iteration
depends on the w computed at the previous one. This imposes overhead
because every iteration must re-read the input file into memory. In
Dryad, the whole job can be expressed as a data flow DAG as shown in
Figure 4a, but the data must still must be reloaded from disk at each
iteration. Reusing the data in memory between iterations in Dryad would
require cyclic data flow.

Spark's execution is shown in Figure 4b. Spark uses the long-lived
nature of Mesos executors to cache a slice of the dataset in memory at
each executor, and then run multiple iterations on this cached data.
This caching is achieved in a fault-tolerant manner: if a node is lost,
Spark remembers how to recompute its slice of the data.

By building Spark on top of Mesos, we were able to keep its
implementation small (about 1300 lines of code), yet still capable of
outperforming Hadoop by 10× for iterative jobs. In particular, using
Mesos's API saved us the time to write a master daemon, slave daemon,
and communication protocols between them for Spark. The main pieces we
had to write were a framework scheduler (which uses delay scheduling for
locality) and user APIs.

\section{6 Evaluation}\label{evaluation}

We evaluated Mesos through a series of experiments on the Amazon Elastic
Compute Cloud (EC2). We begin with a macrobenchmark that evaluates how
the system shares resources between four workloads, and go on to present
a series of smaller experiments designed to evaluate overhead,
decentralized scheduling, our specialized framework (Spark),
scalability, and failure recovery.

\section{6.1 Macrobenchmark}\label{macrobenchmark}

To evaluate the primary goal of Mesos, which is enabling diverse
frameworks to efficiently share a cluster, we ran a

Bin

Job Type

Map Tasks

Reduce Tasks

\# Jobs Run

1

selection

1

NA

38

2

text search

2

NA

18

3

aggregation

10

2

14

4

selection

50

NA

12

5

aggregation

100

10

6

6

selection

200

NA

6

7

text search

400

NA

4

8

join

400

30

2

Table 3: Job types for each bin in our Facebook Hadoop mix.

macrobenchmark consisting of a mix of four workloads:

\begin{itemize}
\tightlist
\item
  A Hadoop instance running a mix of small and large jobs based on the
  workload at Facebook.\\
\item
  A Hadoop instance running a set of large batch jobs.\\
\item
  Spark running a series of machine learning jobs.\\
\item
  Torque running a series of MPI jobs.
\end{itemize}

We compared a scenario where the workloads ran as four frameworks on a
96-node Mesos cluster using fair sharing to a scenario where they were
each given a static partition of the cluster (24 nodes), and measured
job response times and resource utilization in both cases. We used EC2
nodes with 4 CPU cores and 15 GB of RAM.

We begin by describing the four workloads in more detail, and then
present our results.

\section{6.1.1 Macrobenchmark Workloads}\label{macrobenchmark-workloads}

Facebook Hadoop Mix Our Hadoop job mix was based on the distribution of
job sizes and inter-arrival times at Facebook, reported in {[}38{]}. The
workload consists of 100 jobs submitted at fixed times over a 25-minute
period, with a mean inter-arrival time of 14s. Most of the jobs are
small (1-12 tasks), but there are also large jobs of up to 400 tasks.
\(^{4}\) The jobs themselves were from the Hive benchmark {[}6{]}, which
contains four types of queries: text search, a simple selection, an
aggregation, and a join that gets translated into multiple MapReduce
steps. We grouped the jobs into eight bins of job type and size (listed
in Table 3) so that we could compare performance in each bin. We also
set the framework scheduler to perform fair sharing between its jobs, as
this policy is used at Facebook.

Large Hadoop Mix To emulate batch workloads that need to run
continuously, such as web crawling, we had a second instance of Hadoop
run a series of IO-intensive 2400-task text search jobs. A script
launched ten of these jobs, submitting each one after the previous one
finished.

Spark We ran five instances of an iterative machine learning job on
Spark. These were launched by a script that waited 2 minutes after each
job ended to submit the next. The job we used was alternating least
squares

\begin{enumerate}
\def\labelenumi{(\alph{enumi})}
\tightlist
\item
  Facebook Hadoop Mix\\
  \textit{[Image omitted: images/00ae7ae578f6d79aea4fbe2c320bcc2ea396790f26831c551d49b48d3617fec0.jpg]}
\end{enumerate}

line

{\def\LTcaptype{none} % do not increment counter
\begin{longtable}[]{@{}lll@{}}
\toprule\noalign{}
Time (s) & Static Partitioning & Mesos \\
\midrule\noalign{}
\endhead
\bottomrule\noalign{}
\endlastfoot
0 & 1.0 & 0.0 \\
100 & 0.2 & 0.0 \\
200 & 0.3 & 0.0 \\
300 & 0.6 & 0.0 \\
400 & 0.2 & 0.0 \\
500 & 0.4 & 0.0 \\
600 & 0.8 & 0.0 \\
700 & 0.2 & 0.0 \\
800 & 0.4 & 0.0 \\
900 & 0.2 & 0.0 \\
1000 & 0.1 & 0.0 \\
1100 & 0.6 & 0.0 \\
1200 & 0.2 & 0.0 \\
1300 & 0.3 & 0.0 \\
1400 & 0.2 & 0.0 \\
1500 & 0.4 & 0.0 \\
1600 & 0.1 & 0.0 \\
\end{longtable}
}

\begin{enumerate}
\def\labelenumi{(\alph{enumi})}
\setcounter{enumi}{1}
\tightlist
\item
  Large Hadoop Mix\\
  \textit{[Image omitted: images/1c1a3ca98620fb7bbaf7d0e6bce4d9c5c94db7c32960121a43f44a056e20b21b.jpg]}
\end{enumerate}

line

{\def\LTcaptype{none} % do not increment counter
\begin{longtable}[]{@{}lll@{}}
\toprule\noalign{}
Time (s) & Static Partitioning & Mesos \\
\midrule\noalign{}
\endhead
\bottomrule\noalign{}
\endlastfoot
0 & 0.2 & 0.6 \\
500 & 0.2 & 0.8 \\
1000 & 0.2 & 0.9 \\
1500 & 0.2 & 0.8 \\
2000 & 0.2 & 0.2 \\
2500 & 0.2 & 0.2 \\
3000 & 0.2 & 0.2 \\
\end{longtable}
}

\begin{enumerate}
\def\labelenumi{(\alph{enumi})}
\setcounter{enumi}{2}
\tightlist
\item
  Spark\\
  \textit{[Image omitted: images/60ee5a81ddbdb35c9ea7a2309ab67948e31cbcf92e18f90f4eac57534da7a2ad.jpg]}
\end{enumerate}

line

{\def\LTcaptype{none} % do not increment counter
\begin{longtable}[]{@{}lll@{}}
\toprule\noalign{}
Time (s) & Static Partitioning & Mesos \\
\midrule\noalign{}
\endhead
\bottomrule\noalign{}
\endlastfoot
0 & 0.2 & 0.4 \\
100 & 0.3 & 0.8 \\
200 & 0.2 & 0.4 \\
300 & 0.2 & 0.4 \\
400 & 0.2 & 0.4 \\
500 & 0.2 & 0.2 \\
600 & 0.2 & 0.3 \\
700 & 0.2 & 0.4 \\
800 & 0.2 & 0.4 \\
900 & 0.2 & 0.8 \\
1000 & 0.2 & 0.9 \\
1100 & 0.2 & 0.4 \\
1200 & 0.2 & 0.4 \\
1300 & 0.2 & 0.4 \\
1400 & 0.2 & 0.6 \\
1500 & 0.2 & 0.4 \\
1600 & 0.2 & 0.4 \\
1700 & 0.2 & 0.4 \\
\end{longtable}
}

\begin{enumerate}
\def\labelenumi{(\alph{enumi})}
\setcounter{enumi}{3}
\tightlist
\item
  Torque / MPI\\
  \textit{[Image omitted: images/e29eaf707bd769fbea156d7c44902c7a5493487d1457279c11065eda16a1d1f4.jpg]}
\end{enumerate}

bar

{\def\LTcaptype{none} % do not increment counter
\begin{longtable}[]{@{}lll@{}}
\toprule\noalign{}
Time (s) & Static Partitioning & Mesos \\
\midrule\noalign{}
\endhead
\bottomrule\noalign{}
\endlastfoot
0 & 0.0 & 0.0 \\
200 & 0.1 & 0.0 \\
400 & 0.2 & 0.0 \\
600 & 0.3 & 0.2 \\
800 & 0.1 & 0.1 \\
1000 & 0.0 & 0.0 \\
1200 & 0.2 & 0.0 \\
1400 & 0.3 & 0.2 \\
1600 & 0.1 & 0.1 \\
\end{longtable}
}

Figure 5: Comparison of cluster shares (fraction of CPUs) over time for
each of the frameworks in the Mesos and static partitioning
macrobenchmark scenarios. On Mesos, frameworks can scale up when their
demand is high and that of other frameworks is low, and thus finish jobs
faster. Note that the plots' time axes are different (e.g., the large
Hadoop mix takes 3200s with static partitioning).

\textit{[Image omitted: images/017e56a2e67d5acee657dbb78ea6d67f7e20de09c3629f87d3a64989ec97aa15.jpg]}

bar\_stacked

{\def\LTcaptype{none} % do not increment counter
\begin{longtable}[]{@{}
  >{\raggedright\arraybackslash}p{(\linewidth - 8\tabcolsep) * \real{0.1356}}
  >{\raggedright\arraybackslash}p{(\linewidth - 8\tabcolsep) * \real{0.0847}}
  >{\raggedright\arraybackslash}p{(\linewidth - 8\tabcolsep) * \real{0.3051}}
  >{\raggedright\arraybackslash}p{(\linewidth - 8\tabcolsep) * \real{0.2712}}
  >{\raggedright\arraybackslash}p{(\linewidth - 8\tabcolsep) * \real{0.2034}}@{}}
\toprule\noalign{}
\begin{minipage}[b]{\linewidth}\raggedright
Time (s)
\end{minipage} & \begin{minipage}[b]{\linewidth}\raggedright
Spark
\end{minipage} & \begin{minipage}[b]{\linewidth}\raggedright
Facebook Hadoop Mix
\end{minipage} & \begin{minipage}[b]{\linewidth}\raggedright
Large Hadoop Mix
\end{minipage} & \begin{minipage}[b]{\linewidth}\raggedright
Torque / MPI
\end{minipage} \\
\midrule\noalign{}
\endhead
\bottomrule\noalign{}
\endlastfoot
0 & 0.0 & 0.0 & 0.0 & 0.0 \\
200 & 0.0 & 0.0 & 0.0 & 0.0 \\
400 & 0.0 & 0.0 & 0.0 & 0.0 \\
600 & 0.0 & 0.0 & 0.0 & 0.0 \\
800 & 0.0 & 0.0 & 0.0 & 0.0 \\
1000 & 0.0 & 0.0 & 0.0 & 0.0 \\
1200 & 0.0 & 0.0 & 0.0 & 0.0 \\
1400 & 0.0 & 0.0 & 0.0 & 0.0 \\
1600 & 0.0 & 0.0 & 0.0 & 0.0 \\
\end{longtable}
}

Figure 6: Framework shares on Mesos during the macrobenchmark. By
pooling resources, Mesos lets each workload scale up to fill gaps in the
demand of others. In addition, fine-grained sharing allows resources to
be reallocated in tens of seconds.

(ALS), a collaborative filtering algorithm {[}42{]}. This job is
CPU-intensive but also benefits from caching its input data on each
node, and needs to broadcast updated parameters to all nodes running its
tasks on each iteration.

Torque / MPI Our Torque framework ran eight instances of the tachyon
raytracing job {[}35{]} that is part of the SPEC MPI2007 benchmark. Six
of the jobs ran small problem sizes and two ran large ones. Both types
used 24 parallel tasks. We submitted these jobs at fixed times to both
clusters. The tachyon job is CPU-intensive.

\section{6.1.2 Macrobenchmark Results}\label{macrobenchmark-results}

A successful result for Mesos would show two things: that Mesos achieves
higher utilization than static partitioning, and that jobs finish at
least as fast in the shared cluster as they do in their static
partition, and possibly faster due to gaps in the demand of other
frameworks. Our results show both effects, as detailed below.

We show the fraction of CPU cores allocated to each

\textit{[Image omitted: images/6ed69a8fba50877a102a9c8083303ded6997c58dd55ce6ce651cfa340f207f66.jpg]}

line

{\def\LTcaptype{none} % do not increment counter
\begin{longtable}[]{@{}lll@{}}
\toprule\noalign{}
Time (s) & Mesos & Static \\
\midrule\noalign{}
\endhead
\bottomrule\noalign{}
\endlastfoot
0 & 0 & 0 \\
200 & 80 & 60 \\
400 & 60 & 40 \\
600 & 80 & 60 \\
800 & 60 & 40 \\
1000 & 80 & 60 \\
1200 & 60 & 40 \\
1400 & 80 & 60 \\
1600 & 0 & 0 \\
\end{longtable}
}

\textit{[Image omitted: images/aabaf68f24f68663c023619d59ef4acfed716032e40002aaacb0686869c7a546.jpg]}

line

{\def\LTcaptype{none} % do not increment counter
\begin{longtable}[]{@{}lll@{}}
\toprule\noalign{}
Time (s) & Mesos & Static \\
\midrule\noalign{}
\endhead
\bottomrule\noalign{}
\endlastfoot
0 & 10 & 10 \\
200 & 30 & 15 \\
400 & 20 & 10 \\
600 & 40 & 15 \\
800 & 30 & 10 \\
1000 & 20 & 15 \\
1200 & 30 & 10 \\
1400 & 40 & 15 \\
1600 & 20 & 10 \\
\end{longtable}
}

Figure 7: Average CPU and memory utilization over time across all nodes
in the Mesos cluster vs.~static partitioning.

framework by Mesos over time in Figure 6. We see that Mesos enables each
framework to scale up during periods when other frameworks have low
demands, and thus keeps cluster nodes busier. For example, at time 350,
when both Spark and the Facebook Hadoop framework have no running jobs
and Torque is using 1/8 of the cluster, the large-job Hadoop framework
scales up to 7/8 of the cluster. In addition, we see that resources are
reallocated rapidly (e.g., when a Facebook Hadoop job starts around time
360) due to the fine-grained nature of tasks. Finally, higher allocation
of nodes also translates into increased CPU and memory utilization (by
10\% for CPU and 17\% for memory), as shown in Figure 7.

A second question is how much better jobs perform under Mesos than when
using a statically partitioned cluster. We present this data in two
ways. First, Figure 5 compares the resource allocation over time of each
framework in the shared and statically partitioned clusters. Shaded
areas show the allocation in the stat-

Framework

Sum of Exec Times w/ Static Partitioning (s)

Sum of Exec Times with Mesos (s)

Speedup

Facebook Hadoop Mix

7235

6319

1.14

Large Hadoop Mix

3143

1494

2.10

Spark

1684

1338

1.26

Torque / MPI

3210

3352

0.96

Table 4: Aggregate performance of each framework in the macrobenchmark
(sum of running times of all the jobs in the framework). The speedup
column shows the relative gain on Mesos.

ically partitioned cluster, while solid lines show the share on Mesos.
We see that the fine-grained frameworks (Hadoop and Spark) take
advantage of Mesos to scale up beyond 1/4 of the cluster when global
demand allows this, and consequently finish bursts of submitted jobs
faster in Mesos. At the same time, Torque achieves roughly similar
allocations and job durations under Mesos (with some differences
explained later).

Second, Tables 4 and 5 show a breakdown of job performance for each
framework. In Table 4, we compare the aggregate performance of each
framework, defined as the sum of job running times, in the static
partitioning and Mesos scenarios. We see the Hadoop and Spark jobs as a
whole are finishing faster on Mesos, while Torque is slightly slower.
The framework that gains the most is the large-job Hadoop mix, which
almost always has tasks to run and fills in the gaps in demand of the
other frameworks; this framework performs 2x better on Mesos.

Table 5 breaks down the results further by job type. We observe two
notable trends. First, in the Facebook Hadoop mix, the smaller jobs
perform worse on Mesos. This is due to an interaction between the fair
sharing performed by Hadoop (among its jobs) and the fair sharing in
Mesos (among frameworks): During periods of time when Hadoop has more
than 1/4 of the cluster, if any jobs are submitted to the other
frameworks, there is a delay before Hadoop gets a new resource offer
(because any freed up resources go to the framework farthest below its
share), so any small job submitted during this time is delayed for a
long time relative to its length. In contrast, when running alone,
Hadoop can assign resources to the new job as soon as any of its tasks
finishes. This problem with hierarchical fair sharing is also seen in
networks {[}34{]}, and could be mitigated by running the small jobs on a
separate framework or using a different allocation policy (e.g., using
lottery scheduling instead of offering all freed resources to the
framework with the lowest share).

Lastly, Torque is the only framework that performed worse, on average,
on Mesos. The large tachyon jobs took on average 2 minutes longer, while
the small ones took 20s longer. Some of this delay is due to Torque
having to wait to launch 24 tasks on Mesos before starting each job, but
the average time this takes is 12s. We believe that the rest of the
delay is due to stragglers (slow nodes). In our standalone Torque run,
we saw two jobs take about 60s longer to run than the others (Fig. 5d).
We discovered that both of these jobs were using a node that performed
slower on single-node benchmarks than the others (in fact, Linux
reported 40\% lower bogomips on it). Because tachyon hands out equal
amounts of work to each node, it runs as slowly as the slowest node.

Framework

Job Type

Exec Time w/ Static Partitioning (s)

Avg. Speedup on Mesos

Facebook Hadoop Mix

selection (1)

24

0.84

text search (2)

31

0.90

aggregation (3)

82

0.94

selection (4)

65

1.40

aggregation (5)

192

1.26

selection (6)

136

1.71

text search (7)

137

2.14

join (8)

662

1.35

Large Hadoop Mix

text search

314

2.21

Spark

ALS

337

1.36

Torque / MPI

small tachyon

261

0.91

large tachyon

822

0.88

Table 5: Performance of each job type in the macrobenchmark. Bins for
the Facebook Hadoop mix are in parentheses.

\textit{[Image omitted: images/f2d0ed78212e2bd7e97d431d81f8df0e42f65d42dd3a93cd1770486486cd6d83.jpg]}

line

{\def\LTcaptype{none} % do not increment counter
\begin{longtable}[]{@{}lll@{}}
\toprule\noalign{}
Scenario & Data Locality (\%) & Job Running Time (s) \\
\midrule\noalign{}
\endhead
\bottomrule\noalign{}
\endlastfoot
Static partitioning & 20 & 600 \\
Mesos, no delay sched. & 50 & 480 \\
Mesos, 1s delay sched. & 90 & 360 \\
Mesos, 5s delay sched. & 100 & 360 \\
\end{longtable}
}

Figure 8: Data locality and average job durations for 16 Hadoop
instances running on a 93-node cluster using static partitioning, Mesos,
or Mesos with delay scheduling.

\section{6.2 Overhead}\label{overhead}

To measure the overhead Mesos imposes when a single framework uses the
cluster, we ran two benchmarks using MPI and Hadoop on an EC2 cluster
with 50 nodes, each with 2 CPU cores and 6.5 GB RAM. We used the
High-Performance LINPACK {[}15{]} benchmark for MPI and a WordCount job
for Hadoop, and ran each job three times. The MPI job took on average
50.9s without Mesos and 51.8s with Mesos, while the Hadoop job took 160s
without Mesos and 166s with Mesos. In both cases, the overhead of using
Mesos was less than 4\%.

\section{6.3 Data Locality through Delay
Scheduling}\label{data-locality-through-delay-scheduling}

In this experiment, we evaluated how Mesos' resource offer mechanism
enables frameworks to control their tasks' placement, and in particular,
data locality. We ran 16 instances of Hadoop using 93 EC2 nodes, each
with 4 CPU cores and 15 GB RAM. Each node ran a

map-only scan job that searched a 100 GB file spread throughout the
cluster on a shared HDFS file system and outputted 1\% of the records.
We tested four scenarios: giving each Hadoop instance its own 5-6 node
static partition of the cluster (to emulate organizations that use
coarse-grained cluster sharing systems), and running all instances on
Mesos using either no delay scheduling, 1s delay scheduling or 5s delay
scheduling.

Figure 8 shows averaged measurements from the 16 Hadoop instances across
three runs of each scenario. Using static partitioning yields very low
data locality (18\%) because the Hadoop instances are forced to fetch
data from nodes outside their partition. In contrast, running the Hadoop
instances on Mesos improves data locality, even without delay
scheduling, because each Hadoop instance has tasks on more nodes of the
cluster (there are 4 tasks per node), and can therefore access more
blocks locally. Adding a 1-second delay brings locality above 90\%, and
a 5-second delay achieves 95\% locality, which is competitive with
running one Hadoop instance alone on the whole cluster. As expected, job
performance improves with data locality: jobs run 1.7x faster in the 5s
delay scenario than with static partitioning.

\section{6.4 Spark Framework}\label{spark-framework-1}

We evaluated the benefit of running iterative jobs using the specialized
Spark framework we developed on top of Mesos (Section 5.3) over the
general-purpose Hadoop framework. We used a logistic regression job
implemented in Hadoop by machine learning researchers in our lab, and
wrote a second version of the job using Spark. We ran each version
separately on 20 EC2 nodes, each with 4 CPU cores and 15 GB RAM. Each
experiment used a 29 GB data file and varied the number of logistic
regression iterations from 1 to 30 (see Figure 9).

With Hadoop, each iteration takes 127s on average, because it runs as a
separate MapReduce job. In contrast, with Spark, the first iteration
takes 174s, but subsequent iterations only take about 6 seconds, leading
to a speedup of up to 10x for 30 iterations. This happens because the
cost of reading the data from disk and parsing it is much higher than
the cost of evaluating the gradient function computed by the job on each
iteration. Hadoop incurs the read/parsing cost on each iteration, while
Spark reuses cached blocks of parsed data and only incurs this cost
once. The longer time for the first iteration in Spark is due to the use
of slower text parsing routines.

\section{6.5 Mesos Scalability}\label{mesos-scalability}

To evaluate Mesos' scalability, we emulated large clusters by running up
to 50,000 slave daemons on 99 Amazon EC2 nodes, each with 8 CPU cores
and 6 GB RAM. We used one EC2 node for the master and the rest of the
nodes to run slaves. During the experiment, each of 200 frameworks
running throughout the cluster continuously launches tasks, starting one
task on each slave that it receives a resource offer for. Each task
sleeps for a period of time based on a normal distribution with a mean
of 30 seconds and standard deviation of 10s, and then ends. Each slave
runs up to two tasks at a time.

\textit{[Image omitted: images/c4544f5f0c4678bcbec461040f1a624e9d9c9b40e87fea3ae7f4eca9b0ada2fb.jpg]}

line

{\def\LTcaptype{none} % do not increment counter
\begin{longtable}[]{@{}lll@{}}
\toprule\noalign{}
Number of Iterations & Hadoop & Spark \\
\midrule\noalign{}
\endhead
\bottomrule\noalign{}
\endlastfoot
0 & 0 & 0 \\
10 & 1200 & 200 \\
20 & 2600 & 250 \\
30 & 3900 & 300 \\
\end{longtable}
}

Figure 9: Hadoop and Spark logistic regression running times.\\
\textit{[Image omitted: images/1ae84cab2049d038e4e4b471c6e53b083a9f6c11cc3b7faec8c434a98b9a173b.jpg]}

line

{\def\LTcaptype{none} % do not increment counter
\begin{longtable}[]{@{}ll@{}}
\toprule\noalign{}
Number of Nodes & Task Launch Overhead (seconds) \\
\midrule\noalign{}
\endhead
\bottomrule\noalign{}
\endlastfoot
0 & 0.15 \\
10000 & 0.20 \\
20000 & 0.35 \\
30000 & 0.35 \\
40000 & 0.55 \\
50000 & 0.75 \\
\end{longtable}
}

Figure 10: Mesos master's scalability versus number of slaves.

Once the cluster reached steady-state (i.e., the 200 frameworks achieve
their fair shares and all resources were allocated), we launched a test
framework that runs a single 10 second task and measured how long this
framework took to finish. This allowed us to calculate the extra delay
incurred over 10s due to having to register with the master, wait for a
resource offer, accept it, wait for the master to process the response
and launch the task on a slave, and wait for Mesos to report the task as
finished.

We plot this extra delay in Figure 10, showing averages of 5 runs. We
observe that the overhead remains small (less than one second) even at
50,000 nodes. In particular, this overhead is much smaller than the
average task and job lengths in data center workloads (see Section 2).
Because Mesos was also keeping the cluster fully allocated, this
indicates that the master kept up with the load placed on it.
Unfortunately, the EC2 virtualized environment limited scalability
beyond 50,000 slaves, because at 50,000 slaves the master was processing
100,000 packets per second (in+out), which has been shown to be the
current achievable limit on EC2 {[}12{]}.

\section{6.6 Failure Recovery}\label{failure-recovery}

To evaluate recovery from master failures, we conducted an experiment
with 200 to 4000 slave daemons on 62 EC2 nodes with 4 cores and 15 GB
RAM. We ran 200 frameworks that each launched 20-second tasks, and two
Mesos masters connected to a 5-node ZooKeeper quorum. We synchronized
the two masters' clocks using NTP

and measured the mean time to recovery (MTTR) after killing the active
master. The MTTR is the time for all of the slaves and frameworks to
connect to the second master. In all cases, the MTTR was between 4 and 8
seconds, with 95\% confidence intervals of up to 3s on either side.

\section{6.7 Performance Isolation}\label{performance-isolation}

As discussed in Section 3.4, Mesos leverages existing OS isolation
mechanism to provide performance isolation between different frameworks'
tasks running on the same slave. While these mechanisms are not perfect,
a preliminary evaluation of Linux Containers {[}9{]} shows promising
results. In particular, using Containers to isolate CPU usage between a
MediaWiki web server (consisting of multiple Apache processes running
PHP) and a ``hog'' application (consisting of 256 processes spinning in
infinite loops) shows on average only a 30\% increase in request latency
for Apache versus a 550\% increase when running without Containers. We
refer the reader to {[}29{]} for a fuller evaluation of OS isolation
mechanisms.

\section{7 Related Work}\label{related-work}

HPC and Grid Schedulers. The high performance computing (HPC) community
has long been managing clusters {[}33, 41{]}. However, their target
environment typically consists of specialized hardware, such as
Infiniband and SANs, where jobs do not need to be scheduled local to
their data. Furthermore, each job is tightly coupled, often using
barriers or message passing. Thus, each job is monolithic, rather than
composed of fine-grained tasks, and does not change its resource demands
during its lifetime. For these reasons, HPC schedulers use centralized
scheduling, and require users to declare the required resources at job
submission time. Jobs are then given coarse-grained allocations of the
cluster. Unlike the Mesos approach, this does not allow jobs to locally
access data distributed across the cluster. Furthermore, jobs cannot
grow and shrink dynamically. In contrast, Mesos supports fine-grained
sharing at the level of tasks and allows frameworks to control their
placement.

Grid computing has mostly focused on the problem of making diverse
virtual organizations share geographically distributed and separately
administered resources in a secure and interoperable way. Mesos could
well be used within a virtual organization inside a larger grid.

Public and Private Clouds. Virtual machine clouds such as Amazon EC2
{[}1{]} and Eucalyptus {[}31{]} share common goals with Mesos, such as
isolating applications while providing a low-level abstraction (VMs).
However, they differ from Mesos in several important ways. First, their
relatively coarse grained VM allocation model leads to less efficient
resource utilization and data sharing than in Mesos. Second, these
systems generally do not let applications specify placement needs beyond
the size of VM they require. In contrast, Mesos allows frameworks to be
highly selective about task placement.

Quincy. Quincy {[}25{]} is a fair scheduler for Dryad that uses a
centralized scheduling algorithm for Dryad's DAG-based programming
model. In contrast, Mesos provides the lower-level abstraction of
resource offers to support multiple cluster computing frameworks.

Condor. The Condor cluster manager uses the Class-Ads language {[}32{]}
to match nodes to jobs. Using a resource specification language is not
as flexible for frameworks as resource offers, since not all
requirements may be expressible. Also, porting existing frameworks,
which have their own schedulers, to Condor would be more difficult than
porting them to Mesos, where existing schedulers fit naturally into the
two-level scheduling model.

Next-Generation Hadoop. Recently, Yahoo! announced a redesign for Hadoop
that uses a two-level scheduling model, where per-application masters
request resources from a central manager {[}14{]}. The design aims to
support non-MapReduce applications as well. While details about the
scheduling model in this system are currently unavailable, we believe
that the new application masters could naturally run as Mesos
frameworks.

\section{8 Conclusion and Future Work}\label{conclusion-and-future-work}

We have presented Mesos, a thin management layer that allows diverse
cluster computing frameworks to efficiently share resources. Mesos is
built around two design elements: a fine-grained sharing model at the
level of tasks, and a distributed scheduling mechanism called resource
offers that delegates scheduling decisions to the frameworks. Together,
these elements let Mesos achieve high utilization, respond quickly to
workload changes, and cater to diverse frameworks while remaining
scalable and robust. We have shown that existing frameworks can
effectively share resources using Mesos, that Mesos enables the
development of specialized frameworks providing major performance gains,
such as Spark, and that Mesos's simple design allows the system to be
fault tolerant and to scale to 50,000 nodes.

In future work, we plan to further analyze the resource offer model and
determine whether any extensions can improve its efficiency while
retaining its flexibility. In particular, it may be possible to have
frameworks give richer hints about offers they would like to receive.
Nonetheless, we believe that below any hint system, frameworks should
still have the ability to reject offers and to choose which tasks to
launch on each resource, so that their evolution is not constrained by
the hint language provided by the system.

We are also currently using Mesos to manage resources on a 40-node
cluster in our lab and in a test deployment at Twitter, and plan to
report on lessons from

these deployments in future work.

\section{9 Acknowledgements}\label{acknowledgements}

We thank our industry colleagues at Google, Twitter, Facebook, Yahoo!
and Cloudera for their valuable feedback on Mesos. This research was
supported by California MICRO, California Discovery, the Natural
Sciences and Engineering Research Council of Canada, a National Science
Foundation Graduate Research Fellowship, \(^{5}\) the Swedish Research
Council, and the following Berkeley RAD Lab sponsors: Google, Microsoft,
Oracle, Amazon, Cisco, Cloudera, eBay, Facebook, Fujitsu, HP, Intel,
NetApp, SAP, VMware, and Yahoo!.

\section{References}\label{references}

{[}1{]} Amazon EC2. http://aws.amazon.com/ec2.\\
{[}2{]} Apache Hadoop. http://hadoop.apache.org.\\
{[}3{]} Apache Hive. http://hadoop.apache.org/hive.\\
{[}4{]} Apache ZooKeeper. hadoop.apache.org/zookeeper.\\
{[}5{]} Hive -- A Petabyte Scale Data Warehouse using Hadoop.
http://www.facebook.com/note.php?note\_id=89508453919.\\
{[}6{]} Hive performance benchmarks.
http://issues.apache.org/jira/browse/HIVE-396.\\
{[}7{]} LibProcess Homepage. http:
//www.eecs.berkeley.edu/\textasciitilde benh/libprocess.\\
{[}8{]} Linux 2.6.33 release notes.
http://kernelnewbies.org/Linux\_2\_6\_33.\\
{[}9{]} Linux containers (LXC) overview document.
http://lxc.sourceforge.net/lxc.html.\\
{[}10{]} Personal communication with Dhruba Borthakur from Facebook.\\
{[}11{]} Personal communication with Owen O'Malley and Arun C. Murthy
from the Yahoo! Hadoop team.\\
{[}12{]} RightScale blog.
blog.rightscale.com/2010/04/01/benchmarking-load-balancers-in-the-cloud.\\
{[}13{]} Solaris Resource Management.
http://docs.sun.com/app/docs/doc/817-1592.\\
{[}14{]} The Next Generation of Apache Hadoop MapReduce.
http://developer.yahoo.com/blogs/hadoop/posts/2011/02/mapreduce-nextgen.\\
{[}15{]} E. Anderson, Z. Bai, J. Dongarra, A. Greenbaum, A. McKenney, J.
Du Croz, S. Hammerling, J. Demmel, C. Bischof, and D. Sorensen. LAPACK:
a portable linear algebra library for high-performance computers. In
Supercomputing '90, 1990.\\
{[}16{]} A. Bouteiller, F. Cappello, T. Herault, G. Krawezik, P.
Lemarinier, and F. Magniette. Mpich-v2: a fault tolerant MPI for
volatile nodes based on pessimistic sender based message logging. In
Supercomputing '03, 2003.\\
{[}17{]} T. Condie, N. Conway, P. Alvaro, and J. M. Hellerstein.
MapReduce online. In NSDI '10, May 2010.\\
{[}18{]} J. Dean and S. Ghemawat. MapReduce: Simplified data processing
on large clusters. In OSDI, pages 137--150, 2004.\\
{[}19{]} J. Ekanayake, H. Li, B. Zhang, T. Gunarathne, S.-H. Bae, J.
Qiu, and G. Fox. Twister: a runtime for iterative mapreduce. In Proc.
HPDC '10, 2010.\\
{[}20{]} D. R. Engler, M. F. Kaashoek, and J. O'Toole. Exokernel: An
operating system architecture for application-level resource management.
In SOSP, pages 251--266, 1995.\\
{[}21{]} A. Ghodsi, M. Zaharia, B. Hindman, A. Konwinski, S. Shenker,
and I. Stoica. Dominant resource fairness: fair allocation of multiple
resource types. In NSDI, 2011.

{[}22{]} T. Hastie, R. Tibshirani, and J. Friedman. The Elements of
Statistical Learning: Data Mining, Inference, and Prediction. Springer
Publishing Company, New York, NY, 2009.\\
{[}23{]} B. Hindman, A. Konwinski, M. Zaharia, A. Ghodsi, A. D. Joseph,
R. H. Katz, S. Shenker, and I. Stoica. Mesos: A platform for
fine-grained resource sharing in the data center. Technical Report
UCB/EECS-2010-87, UC Berkeley, May 2010.\\
{[}24{]} M. Isard, M. Budiu, Y. Yu, A. Birrell, and D. Fetterly. Dryad:
distributed data-parallel programs from sequential building blocks. In
EuroSys 07, 2007.\\
{[}25{]} M. Isard, V. Prabhakaran, J. Currey, U. Wieder, K. Talwar, and
A. Goldberg. Quincy: Fair scheduling for distributed computing clusters.
In SOSP, November 2009.\\
{[}26{]} S. Y. Ko, I. Hoque, B. Cho, and I. Gupta. On availability of
intermediate data in cloud computations. In HOTOS, May 2009.\\
{[}27{]} D. Logothetis, C. Olston, B. Reed, K. C. Webb, and K. Yocum.
Stateful bulk processing for incremental analytics. In Proc. ACM
symposium on Cloud computing, SoCC '10, 2010.\\
{[}28{]} G. Malewicz, M. H. Austern, A. J. Bik, J. C. Dehnert, I. Horn,
N. Leiser, and G. Czajkowski. Pregel: a system for large-scale graph
processing. In SIGMOD, pages 135--146, 2010.\\
{[}29{]} J. N. Matthews, W. Hu, M. Hapuarachchi, T. Deshane, D. Dimatos,
G. Hamilton, M. McCabe, and J. Owens. Quantifying the performance
isolation properties of virtualization systems. In ExpCS '07, 2007.\\
{[}30{]} D. G. Murray, M. Schwarzkopf, C. Smowton, S. Smith, A.
Madhavapeddy, and S. Hand. Ciel: a universal execution engine for
distributed data-flow computing. In NSDI, 2011.\\
{[}31{]} D. Nurmi, R. Wolski, C. Grzegorczyk, G. Obertelli, S. Soman, L.
Youseff, and D. Zagorodnov. The Eucalyptus open-source cloud-computing
system. In CCA '08, 2008.\\
{[}32{]} R. Raman, M. Livny, and M. Solomon. Matchmaking: An extensible
framework for distributed resource management. Cluster Computing,
2:129--138, April 1999.\\
{[}33{]} G. Staples. TORQUE resource manager. In Proc. Supercomputing
'06, 2006.\\
{[}34{]} I. Stoica, H. Zhang, and T. S. E. Ng. A hierarchical fair
service curve algorithm for link-sharing, real-time and priority
services. In SIGCOMM '97, pages 249--262, 1997.\\
{[}35{]} J. Stone. Tachyon ray tracing system.
http://jedi.ks.uiuc.edu/\textasciitilde johns/raytracer.\\
{[}36{]} C. A. Waldspurger and W. E. Weihl. Lottery scheduling: flexible
proportional-share resource management. In OSDI, 1994.\\
{[}37{]} Y. Yu, P. K. Gunda, and M. Isard. Distributed aggregation for
data-parallel computing: interfaces and implementations. In SOSP '09,
pages 247--260, 2009.\\
{[}38{]} M. Zaharia, D. Borthakur, J. Sen Sarma, K. Elmeleegy, S.
Shenker, and I. Stoica. Delay scheduling: A simple technique for
achieving locality and fairness in cluster scheduling. In EuroSys 10,
2010.\\
{[}39{]} M. Zaharia, M. Chowdhury, M. J. Franklin, S. Shenker, and I.
Stoica. Spark: cluster computing with working sets. In Proc. HotCloud
'10, 2010.\\
{[}40{]} M. Zaharia, A. Konwinski, A. D. Joseph, R. Katz, and I. Stoica.
Improving MapReduce performance in heterogeneous environments. In Proc.
OSDI '08, 2008.\\
{[}41{]} S. Zhou. LSF: Load sharing in large-scale heterogeneous
distributed systems. In Workshop on Cluster Computing, 1992.\\
{[}42{]} Y. Zhou, D. Wilkinson, R. Schreiber, and R. Pan. Large-scale
parallel collaborative filtering for the Netflix prize. In AAIM, pages
337--348. Springer-Verlag, 2008.

\end{document}
